\chapter{کلیات پژوهش}

\section{مقدمه}
سامانه‌های نرم‌افزاری امروزی معمولاً به صورت پیوسته، توزیع‌شده و تحت بارهای متغیر اجرا می‌شوند. در چنین محیطی، خطاهای عملکردی الزاماً به شکل توقف کامل برنامه دیده نمی‌شوند؛ گاهی افزایش زمان پاسخ، مصرف غیرعادی پردازنده، فشار حافظه، ازدحام قفل‌ها یا الگوی غیرمنتظره ورودی/خروجی فایل باعث افت کیفیت سرویس می‌شود. بنابراین، پایش عملکرد باید بتواند رخدادهای درون برنامه و نشانه‌های سطح سیستم‌عامل را به صورت جریان داده دریافت کرده و به سرعت به معیارهای قابل تحلیل تبدیل کند.

پروژه مبنای این پایان‌نامه، با نام \lr{CE-FINAL-PROJECT}، یک خط لوله پایش و تشخیص ناهنجاری برای برنامه‌های جاوا است. این سامانه از \lr{Java Flight Recorder} یا \lr{JFR} برای مشاهده رخدادهای درون ماشین مجازی جاوا، از \lr{eBPF} برای جمع‌آوری رخدادهای سطح هسته، از \lr{Kafka} برای انتقال پیام، از \lr{Kafka Streams} برای تجمیع پنجره‌ای، و از یک سرویس پایتونی برای تحلیل و گزارش ناهنجاری استفاده می‌کند \cite{cefinalproject,jfrdocs,kafkastreams,ebpfdocs}.

\section{بیان مسئله}
ابزارهای پایش سنتی اغلب یا به معیارهای کلی سیستم‌عامل محدود می‌شوند یا رخدادهای سطح برنامه را بدون ارتباط با زمینه اجرایی سیستم تحلیل می‌کنند. این جدایی باعث می‌شود ریشه‌یابی رخدادهایی مانند افزایش مصرف پردازنده، فشار حافظه، کندی فایل، یا ازدحام قفل دشوار شود. مسئله اصلی این پژوهش، طراحی خط لوله‌ای است که بتواند رخدادهای سطح \lr{JVM} و رخدادهای سطح هسته را به صورت یک جریان یکپارچه دریافت کند، آن‌ها را در بازه‌های زمانی مشخص تجمیع کند و سپس بر اساس قواعد قابل توسعه، ناهنجاری‌های عملکردی را تشخیص دهد.

\section{اهمیت و ضرورت}
پایش عملکرد در محیط تولید باید چند ویژگی کلیدی داشته باشد. نخست، دریافت رخداد باید با سربار قابل قبول انجام شود تا خود ابزار پایش منبع اختلال نباشد. دوم، داده خام باید به معیارهای خلاصه و قابل مقایسه تبدیل شود. سوم، معماری باید توسعه‌پذیر باشد تا اضافه کردن نوع رخداد، معیار یا قانون تشخیص جدید نیازمند بازنویسی کل سامانه نباشد. چهارم، خروجی سامانه باید قابل گزارش‌گیری و ارزیابی باشد. پروژه بررسی‌شده این ویژگی‌ها را با جداسازی اجزای دریافت، تجمیع، تحلیل و ارزیابی دنبال می‌کند.

\section{اهداف پژوهش}
اهداف اصلی این پایان‌نامه عبارت‌اند از:
\begin{itemize}
\item بررسی نیازمندی‌های یک سامانه پایش عملکرد برای برنامه‌های جاوا؛
\item طراحی معماری خط لوله‌ای برای دریافت رخدادهای \lr{JFR} و رخدادهای سطح هسته؛
\item پیاده‌سازی تجمیع پنجره‌ای رخدادها با \lr{Kafka Streams}؛
\item پیاده‌سازی آشکارسازهای ناهنجاری برای پردازنده، حافظه، فایل، شبکه، قفل و رخدادهای هسته؛
\item طراحی سناریوهای ارزیابی کنترل‌شده برای سنجش توان تشخیص سامانه.
\end{itemize}

\section{پرسش‌های پژوهش}
این پژوهش به پرسش‌های زیر پاسخ می‌دهد:
\begin{enumerate}
\item چگونه می‌توان رخدادهای \lr{JFR} و رخدادهای سطح هسته را در یک خط لوله مشترک پردازش کرد؟
\item چه ساختاری برای تجمیع رخدادهای عملکردی در پنجره‌های زمانی مناسب است؟
\item چگونه می‌توان آشکارسازهای ناهنجاری را به صورت مستقل و قابل توسعه طراحی کرد؟
\item چه سناریوهایی برای ارزیابی رفتار سامانه پایش مناسب هستند؟
\end{enumerate}

\section{ساختار پایان‌نامه}
در فصل دوم مفاهیم و فناوری‌های پایه مانند \lr{JFR}، \lr{eBPF}، \lr{Kafka} و تشخیص ناهنجاری معرفی می‌شود. فصل سوم معماری سامانه و جریان داده را توضیح می‌دهد. فصل چهارم جزئیات پیاده‌سازی اجزای اصلی پروژه را بررسی می‌کند. فصل پنجم روش ارزیابی، جمع‌بندی و پیشنهادهای توسعه آینده را ارائه می‌دهد.
